% =============================================================================
% Chapter 12: ML Operations at Scale — Vol II Lecture Slides
% =============================================================================
% IMAGES:
%   - n-models-complexity.pdf    - platform-roi.pdf
%   - mlops-maturity.pdf         - model-registry-deps.pdf
%   - ml-cicd-pipeline.pdf       - staged-rollout.pdf
%   - monitoring-pyramid.pdf     - feature-freshness.pdf
% =============================================================================
\documentclass[aspectratio=169, 12pt]{beamer}
\usepackage{../../assets/beamerthememlsys}

\mlsyssetup{
  volume       = {Volume II},
  chapter      = {Chapter 12},
  logo         = {../../assets/img/logo-mlsysbook.png},
  instlogo     = {../../assets/img/logo-harvard.png},
  chaptertitle = {ML Operations at Scale},
}

% --- Fonts ---
\usepackage{fontspec}
\setsansfont{Helvetica Neue}[
  BoldFont={Helvetica Neue Bold},
  ItalicFont={Helvetica Neue Italic},
  BoldItalicFont={Helvetica Neue Bold Italic},
]
\setmonofont{JetBrains Mono}[Scale=0.85]

% --- Packages ---
\usepackage{booktabs}
\usepackage{amsmath}

% --- Image paths ---
\graphicspath{
  {images/}
}

% --- Chapter-specific macros ---
\newcommand{\TCO}{\text{TCO}_{\text{ML}}}

% --- Helper: safe image include ---
\newcommand{\safeimg}[2][width=\textwidth,keepaspectratio]{%
  \IfFileExists{images/#2}{\includegraphics[#1]{#2}}{%
    \IfFileExists{#2}{\includegraphics[#1]{#2}}{%
      \fbox{\parbox[c][2.5cm][c]{0.85\linewidth}{%
        \centering\footnotesize\textcolor{midgray}{[Missing image]}}}%
    }%
  }%
}

% --- Section count ---
\setcounter{mlsystotalsections}{6}

\title{ML Operations at Scale}
\author{Vijay Janapa Reddi}
\institute{Harvard University}
\date{}

\begin{document}

% =============================================================================
% TITLE SLIDE
% =============================================================================
\mlsystitle{ML Operations at Scale}{From Single-Model to Platform Engineering}{cover_ops_scale.png}

% =============================================================================
% LEARNING OBJECTIVES
% =============================================================================
\begin{frame}{Learning Objectives}
\note{
% -- LINK: Connect to prior chapters
Students built serving infrastructure in Part III. This chapter asks:
what happens when you manage not one model, but a hundred?

% -- NARRATE: What to SAY
Read each objective aloud, pausing on ``platform ROI'' and ``TCO framework.''
These two anchor the quantitative reasoning for the entire chapter.

% -- ENGAGE: Specific question
Ask: ``How many of you have deployed more than one model to production?
At what count did ad hoc practices start breaking?''
Give 10 seconds for a show of hands.

% -- WARN: Specific misconception
Students assume operations scale linearly with model count.
Correct: dependencies grow as O(N^2), alerts as O(N*M).

% -- FLEX: [CORE]
[CORE] Never skip --- objectives frame the entire lecture.
IF AHEAD: Ask students to rank which objective they most want to master.
IF SHORT: Read objectives without elaboration, move on.
}

\small
\begin{enumerate}
  \item Calculate \textbf{platform ROI} to justify shared ML infrastructure
  \item Design \textbf{dependency-aware model registries} for multi-model portfolios
  \item Implement \textbf{CI/CD pipelines} tailored to model risk profiles
  \item Architect \textbf{hierarchical monitoring} that prevents alert fatigue
  \item Quantify \textbf{ML technical debt} using deployment velocity and toil metrics
  \item Compare \textbf{centralized}, \textbf{embedded}, and \textbf{hybrid} organizational patterns
  \item Apply the \textbf{TCO framework} to make infrastructure investment decisions
\end{enumerate}

\end{frame}

\begin{frame}{Visual Language}
\note{
% -- NARRATE: What to SAY
Point to each card: ``Blue = compute, green = data, orange = routing,
red = error. These are consistent across every diagram in this course.
When you see red in a pipeline diagram, something is bottlenecked.''

% -- FLEX: [OPTIONAL]
[OPTIONAL] Skip if students have seen this in a previous chapter deck.
IF SHORT: Say ``same color system as last lecture'' and advance.
}

\small
Throughout this course, colors carry meaning:

\vspace{0.3cm}
\begin{columns}[T]
  \begin{column}{0.45\textwidth}
    \begin{mlsyscard}{computestroke}
    \textbf{Blue} --- Compute / Processing\\
    {\footnotesize GPU ops, forward/backward pass, inference}
    \end{mlsyscard}
    \vspace{0.15cm}
    \begin{mlsyscard}{datastroke}
    \textbf{Green} --- Data / Memory\\
    {\footnotesize Data flow, caches, healthy paths}
    \end{mlsyscard}
  \end{column}
  \begin{column}{0.45\textwidth}
    \begin{mlsyscard}{routingstroke}
    \textbf{Orange} --- Routing / Scheduling\\
    {\footnotesize Load balancers, batch windows}
    \end{mlsyscard}
    \vspace{0.15cm}
    \begin{mlsyscard}{errorstroke}
    \textbf{Red} --- Error / Cost / Bottleneck\\
    {\footnotesize Loss, decode phase, waste}
    \end{mlsyscard}
  \end{column}
\end{columns}
\end{frame}



% =============================================================================
\section{Platform Operations}
% =============================================================================

\begin{frame}{The N-Models Problem}
\note{
% -- LINK: Connect to prior concept
Students just saw the learning objectives listing platform ROI and dependency
management. This slide makes the problem visceral: why platforms exist.

% -- NARRATE: What to SAY
Point to the diagram: ``At 10 models, a few shared data sources. At 50,
the dependency graph is a hairball. At 100, a single upstream change
cascades unpredictably.'' Trace the quadratic growth curve.

% -- ENGAGE: Specific question
Ask: ``At your organization, how many models share the same data sources?''
Expected answer: most students underestimate --- typical is 5--10 shared sources.

% -- WARN: Specific misconception
Students assume managing 100 models is 100x the work of managing one.
Correct: dependencies grow as O(N^2), so 100 models is closer to 10,000x
the coordination complexity.

% -- FLEX: [CORE]
[CORE] This motivates the entire chapter.
IF AHEAD: Ask ``At what N does your dependency graph become unmanageable?''
IF SHORT: Show diagram, state O(N^2), move on.
}

% --- Full-width diagram ---
\centering
\safeimg[width=0.92\textwidth,height=4.5cm,keepaspectratio]{n-models-complexity.pdf}

\vspace{0.1cm}
\mlsysinsight{Key:}{Per-model practices collapse at $\sim$50 models. Dependencies grow as $O(N^2)$.}

\end{frame}

\begin{frame}{Operational Complexity Growth}
\note{
% -- LINK: Connect to prior concept
The N-models diagram showed the complexity curve. This table puts
concrete operational labels on each step of that curve.

% -- NARRATE: What to SAY
Walk column by column: ``At 1 model, monitoring is a single dashboard.
At 100, you have 100 dashboards nobody reads. Debugging shifts from
local to distributed tracing --- a qualitatively different skill.''

% -- ENGAGE: Specific question
Ask: ``Which column transitions from manageable to unmanageable first?''
Expected answer: monitoring (it is the first to break because alert
volume grows as N times M metrics).

% -- WARN: Specific misconception
Students focus on deployment coordination but miss that monitoring
breaks first. Alert fatigue precedes deployment chaos.

% -- FLEX: [OPTIONAL]
[OPTIONAL] This table reinforces the N-models diagram.
IF SHORT: Point to the monitoring row, state the key insight, advance.
IF AHEAD: Ask students to fill in a ``1000 models'' column mentally.
}

\footnotesize
\renewcommand{\arraystretch}{1.15}
\begin{tabular}{@{}llll@{}}
  \toprule
  \textbf{Aspect} & \textbf{1 Model} & \textbf{10 Models} & \textbf{100 Models} \\
  \midrule
  \textbf{Deployment}  & None             & Ad hoc            & Critical path \\
  \textbf{Dependencies}& None             & Some overlap      & Dense graph \\
  \textbf{Monitoring}  & 1 dashboard      & 10 dashboards     & Unmanageable \\
  \textbf{On-call}     & Single team      & Multi-team        & Org-wide \\
  \textbf{Utilization} & Often idle       & Moderate sharing  & Efficiency critical \\
  \textbf{Debugging}   & Local            & Cross-team        & Distributed tracing \\
  \bottomrule
\end{tabular}

\vspace{0.15cm}
\mlsysalert{Consequence:}{100 models with independent failure modes create a constant stream of alerts that masks genuine emergencies.}

\end{frame}

\begin{frame}{Quantifying Platform ROI}
\note{
% -- LINK: Connect to prior concept
The complexity table showed operations becoming unmanageable. This slide
answers: what is the economic case for investing in a platform?

% -- NARRATE: What to SAY
Point to the equation: ``N is model count, T-saved is hours per model,
C-eng is engineer cost. The numerator grows linearly with N; the
denominator is fixed.'' Walk through the worked example: ``50 models,
40 hours each, \$150/hr = \$3.6M before, \$1.6M after. 56\% savings.''

% -- ENGAGE: Specific question
Ask: ``At what model count does your organization need a platform team?''
Give 10 seconds. Expected: most say 50--100; the surprise is 20--50.

% -- WARN: Specific misconception
Students think platform ROI is linear. Correct: it is superlinear because
shared infrastructure amortizes over all models simultaneously.

% -- FLEX: [CORE]
[CORE] The ROI equation is the quantitative anchor for the chapter.
IF AHEAD: Ask ``What if T-saved is only 10 hours instead of 30?''
IF SHORT: Show equation, state the 56\% number, skip the worked example detail.
}

\small
\begin{columns}[T]
  \begin{column}{0.48\textwidth}
    \safeimg[width=\textwidth,height=5.5cm,keepaspectratio]{platform-roi.pdf}
  \end{column}
  \begin{column}{0.48\textwidth}
    $$\text{ROI}_{\text{platform}} = \frac{N \times T_{saved} \times C_{eng}}{C_{platform}}$$

    \vspace{0.15cm}
    {\footnotesize
    \textbf{Worked example} (50 models):
    \begin{itemize}\setlength\itemsep{0pt}
      \item Before: 50 $\times$ 40 hrs $\times$ \$150 = \$3.6M/yr
      \item After: 50 $\times$ 10 hrs $\times$ \$150 + \$667K = \$1.6M/yr
      \item \textcolor{datastroke}{\textbf{Savings: 56\%}} (\$2M/yr)
    \end{itemize}
    }

    \vspace{0.1cm}
    \mlsysconcept{Threshold:}{Break-even at 20--50 models.}
  \end{column}
\end{columns}

\end{frame}

% --- ACTIVE LEARNING: Micro-Retrieval Cue ---
\begin{frame}{Quick Check}
\note{
% -- NARRATE: What to SAY
Pause 15 seconds. Then cold-call. Answer: 20--50 models.
Say: ``Most students guess 100+. The surprise is the threshold is
much lower because shared infrastructure amortizes across all models.''

% -- FLEX: [OPTIONAL]
[OPTIONAL] Micro-retrieval cue reinforcing the ROI slide.
IF SHORT: Skip entirely --- the predict exercise covers this ground.
}

\centering
\vspace{1.0cm}
{\Large\bfseries Quick Check}

\vspace{0.6cm}
{\large At what model count does platform investment\\typically break even?}

\vspace{0.5cm}
{\normalsize\textcolor{midgray}{15 seconds --- then I will cold-call.}}

\end{frame}



\begin{frame}{MLOps Maturity Hierarchy}
\note{
% -- LINK: Connect to prior concept
The ROI equation showed that platforms pay for themselves. This slide
asks: what does the maturity journey look like?

% -- NARRATE: What to SAY
Point to each level: ``L0 = scripts on laptops. L1 = per-model CI/CD,
the most common state. L2 = shared platform, where the superlinear
returns kick in. L3 = enterprise governance across the org.'' Emphasize
the L1-to-L2 transition: ``This is where most organizations stall.''

% -- ENGAGE: Specific question
Ask: ``Where is your organization on this staircase?''
Show of hands for each level. Most will cluster at L0--L1.

% -- WARN: Specific misconception
Students think the jump from L0 to L1 is the hard part. Correct:
L0-to-L1 is just adding CI/CD. The L1-to-L2 jump requires
organizational change --- shared infrastructure, common APIs, platform team.

% -- FLEX: [OPTIONAL]
[OPTIONAL] Reinforces the platform investment argument.
IF SHORT: Show staircase, name the four levels, emphasize L1-to-L2, advance.
}

% --- Full-width diagram ---
\centering
\safeimg[width=0.88\textwidth,height=4.5cm,keepaspectratio]{mlops-maturity.pdf}

\vspace{0.1cm}
\small
\begin{columns}[T]
  \begin{column}{0.22\textwidth}
    \centering
    \textcolor{errorstroke}{\textbf{L0: Manual}}\\
    {\scriptsize Ad hoc, 1--2 models}
  \end{column}
  \begin{column}{0.22\textwidth}
    \centering
    \textcolor{routingstroke}{\textbf{L1: Per-Model}}\\
    {\scriptsize CI/CD per model}
  \end{column}
  \begin{column}{0.22\textwidth}
    \centering
    \textcolor{computestroke}{\textbf{L2: Platform}}\\
    {\scriptsize Shared infra}
  \end{column}
  \begin{column}{0.22\textwidth}
    \centering
    \textcolor{datastroke}{\textbf{L3: Enterprise}}\\
    {\scriptsize Org-wide governance}
  \end{column}
\end{columns}

\end{frame}

% --- ACTIVE LEARNING 1: Predict ---
\begin{frame}{Predict: When Does Platform Investment Pay Off?}
\note{[2 min] Prediction exercise. Students write for 60 seconds.
Do NOT reveal the answer yet. This primes the platform economics discussion.
Ask 2--3 students to share.}

\centering
\vspace{1.0cm}
{\Large\bfseries Think--Write--Share}

\vspace{0.6cm}
{\large Your organization has 30 models, each costing\\
40 engineer-hours/month in operational toil.\\[0.3cm]
\alert{At what model count does a \$2M platform pay for itself?}}

\vspace{0.6cm}
{\normalsize Write your estimate. \textcolor{midgray}{(60 seconds)}}

\end{frame}

% =============================================================================
\section{Multi-Model Management}
% =============================================================================

\begin{frame}{Dependency-Aware Model Registry}
\note{[3 min] The core challenge: upstream changes cascade to downstream
consumers. Walk through the dependency graph. Ask: ``What happens when
the embedding model updates silently?'' Answer: garbage predictions for
all downstream models.}

% --- Full-width diagram ---
\centering
\safeimg[width=0.88\textwidth,height=4.5cm,keepaspectratio]{model-registry-deps.pdf}

\vspace{0.1cm}
\mlsysalert{Without tracking:}{Organizations discover dependencies through production failures.}

\end{frame}

\begin{frame}{Ensemble Management}
\note{[3 min] Recommendation systems invoke 10--50 models per request.
Updating any component requires understanding interactions with all others.
Key patterns: compensation effects, distribution shift propagation, feedback
loops. If short: focus on the staged deployment table.}

\footnotesize
\textbf{A single recommendation request invokes 10--50 models:}

\vspace{0.1cm}
\renewcommand{\arraystretch}{1.1}
\begin{tabular}{@{}llrl@{}}
  \toprule
  \textbf{Deploy Stage} & \textbf{Actions} & \textbf{Duration} & \textbf{Rollback Trigger} \\
  \midrule
  \textbf{Shadow}       & Log only, not served        & 24--48 h  & Quality below threshold \\
  \textbf{Canary}       & 1\% traffic                 & 4--8 h    & Statistical regression \\
  \textbf{Staged}       & 5\% $\to$ 25\% $\to$ 100\% & 24--72 h  & Business metric drop \\
  \textbf{Soak}         & Full traffic, extended watch & 7--14 d   & Delayed effects \\
  \bottomrule
\end{tabular}

\vspace{0.15cm}
\begin{mlsyscard}{routingstroke}
{\footnotesize\textbf{Interaction effects:} A component change improving local metrics might degrade system-level performance through subtle interactions.}
\end{mlsyscard}

\end{frame}

\begin{frame}{Model-Type Operations Diversity}
\note{[2 min] Different model types need fundamentally different operational
patterns. Fraud detection needs hourly updates; LLMs need monthly staged
rollouts. The principle: risk profile determines operational cadence.}

\footnotesize
\renewcommand{\arraystretch}{1.15}
\begin{tabular}{@{}lllll@{}}
  \toprule
  \textbf{Model Type} & \textbf{Update Freq.} & \textbf{Deploy} & \textbf{Primary Risk} & \textbf{Rollback} \\
  \midrule
  \textbf{LLM}   & Monthly     & Staged, careful  & Quality, safety     & Hours--days \\
  \textbf{RecSys} & Daily--weekly & Shadow, A/B     & Engagement drop     & Minutes \\
  \textbf{Fraud}  & Hourly--daily & Rapid + rollback & False negatives     & Seconds \\
  \textbf{Vision} & Weekly--monthly & Canary         & Accuracy regression & Minutes \\
  \bottomrule
\end{tabular}

\vspace{0.15cm}
\mlsysconcept{Principle:}{Risk profile determines operational cadence.}

\vspace{0.1cm}
{\footnotesize LLMs operate slowly because quality regressions are hard to detect. Fraud detection operates continuously because adversaries do not wait.}

\end{frame}

% =============================================================================
\section{CI/CD at Scale}
% =============================================================================

\begin{frame}{ML CI/CD Pipeline}
\note{[3 min] Walk through the pipeline stages and gates. Key difference
from software CI/CD: must validate data, model performance, and fairness ---
not just code correctness. Each stage is idempotent. Ask: ``Which gate
is most likely to catch a silent regression?''}

% --- Full-width diagram ---
\centering
\safeimg[width=0.92\textwidth,height=4.8cm,keepaspectratio]{ml-cicd-pipeline.pdf}

\vspace{0.1cm}
\mlsysinsight{Key difference:}{ML CI/CD validates \textbf{data + model + fairness}, not just code.}

\end{frame}

\begin{frame}{Staged Rollout Strategy}
\note{[3 min] Walk through the four-phase rollout. Emphasize the 20$\times$
risk reduction from 5\% canary. The canary duration formula connects sample
size to traffic percentage. If short: focus on the risk reduction number.}

% --- Full-width diagram ---
\centering
\safeimg[width=0.88\textwidth,height=4.0cm,keepaspectratio]{staged-rollout.pdf}

\vspace{0.1cm}
\small
\textbf{Canary duration:} $t_{stage} = \frac{n_{samples}}{r_{requests} \times p_{stage}}$

\vspace{0.1cm}
{\footnotesize At 1M req/hr with 1\% canary: 10K samples needed $\to$ 1 hour minimum.\\
Total confident rollout: $\sim$4 hours (vs.\ instant blue-green with no safety net).}

\end{frame}

% --- ACTIVE LEARNING 2: Exercise ---
\begin{frame}{Your Turn: Cost of a Delayed Alert}
\note{[3 min] Give students 90 seconds. Answer: 432M requests $\times$ 0.005
lost CTR $\times$ \$0.50/click = \$1.08M. The insight: a ``minor'' 0.5\%
regression at scale is a million-dollar daily disaster.}

\small
\begin{columns}[T]
  \begin{column}{0.55\textwidth}
    {\normalsize\bfseries Calculate the Loss}

    \vspace{0.2cm}
    A recommendation model has a silent bug reducing CTR by 0.5\% (5.0\% $\to$ 4.5\%).
    \begin{itemize}
      \item 5,000 QPS, each click worth \$0.50
      \item Detection takes 24 hours
    \end{itemize}

    \vspace{0.15cm}
    \textbf{How much revenue is lost?}

    {\footnotesize\textcolor{midgray}{(90 seconds --- then compare with a neighbor)}}
  \end{column}
  \begin{column}{0.42\textwidth}
    \pause
    \begin{mlsyscard}{errorstroke}
    \textbf{Solution:}\\[0.1cm]
    {\footnotesize
    Requests: 5K $\times$ 86,400 = 432M\\
    Lost clicks: 432M $\times$ 0.005 = 2.16M\\
    Loss: 2.16M $\times$ \$0.50\\[0.1cm]
    $=$ \textcolor{errorstroke}{\textbf{\$1,080,000}}\\[0.1cm]
    Every hour of faster detection saves \textbf{\$45K}.
    }
    \end{mlsyscard}
  \end{column}
\end{columns}

\end{frame}

\begin{frame}{CI/CD Patterns by Model Type}
\note{[2 min] Connect back to model-type diversity. LLMs: quality-gated,
days to weeks. Fraud: threshold-gated, hours with seconds rollback.
The key: one-size-fits-all CI/CD fails at platform scale.}

\footnotesize
\renewcommand{\arraystretch}{1.15}
\begin{tabular}{@{}lllll@{}}
  \toprule
  \textbf{Pattern} & \textbf{Model Type} & \textbf{Validation} & \textbf{Rollout} & \textbf{Rollback} \\
  \midrule
  \textbf{Quality-gated}   & LLM   & Human eval + safety  & Days--weeks & Hours \\
  \textbf{Metric-driven}   & RecSys & Engagement metrics  & Hours--days & Minutes \\
  \textbf{Threshold-gated} & Fraud  & Precision/recall    & Hours       & Seconds \\
  \textbf{Accuracy-focused} & Vision & Classification     & Days        & Minutes \\
  \bottomrule
\end{tabular}

\vspace{0.15cm}
\begin{mlsyscard}{crimson}
{\footnotesize \textbf{Rollout risk:} $R = P_{\text{regression}} \times I_{\text{regression}} \times E_{\text{exposure}}$. Reduce any factor to reduce total risk.}
\end{mlsyscard}

\end{frame}

% =============================================================================
\section{Monitoring at Scale}
% =============================================================================

\mlsysfocus{The False Alarm Tax}{%
$P(\text{false alert}) = 1 - (1-\alpha)^N$\\[0.3cm]
{\normalsize With $\alpha=0.05$, $N=1000$: $P \approx 1.0$}%
}

\begin{frame}{Alert Fatigue: The Math}
\note{[3 min] This is the key quantitative insight. Even at 3-sigma
specificity, 100 models with 10 metrics generate 864 false alerts/day.
At 2-sigma: 14,400/day. Operators learn to ignore everything.
Ask: ``How many alerts per day before your on-call team gives up?''}

\small
\textbf{100 models $\times$ 10 metrics $\times$ 5-min checks:}

\vspace{0.15cm}
{\footnotesize
\renewcommand{\arraystretch}{1.2}
\begin{tabular}{@{}lrrr@{}}
  \toprule
  \textbf{Threshold} & \textbf{FPR} & \textbf{Daily False Alerts} & \textbf{Status} \\
  \midrule
  2-sigma (95\%)   & 5.0\%  & 14,400 & \textcolor{errorstroke}{\textbf{Useless}} \\
  3-sigma (99.7\%) & 0.3\%  & 864    & \textcolor{routingstroke}{\textbf{Overwhelming}} \\
  4-sigma (99.99\%) & 0.01\% & 29    & \textcolor{datastroke}{\textbf{Manageable}} \\
  \bottomrule
\end{tabular}
}

\vspace{0.15cm}
\mlsysalert{Destructive dynamic:}{Operators ignore alerts $\to$ genuine issues get lost $\to$ monitoring provides \emph{negative} value.}

\end{frame}

\begin{frame}{Hierarchical Monitoring Architecture}
\note{[3 min] The solution to alert fatigue. Four levels: business,
portfolio, model, infrastructure. High-level metrics trigger alarms;
lower-level metrics serve investigation only. Walk through the pyramid.}

\small
\begin{columns}[T]
  \begin{column}{0.48\textwidth}
    \safeimg[width=\textwidth,height=5.5cm,keepaspectratio]{monitoring-pyramid.pdf}
  \end{column}
  \begin{column}{0.48\textwidth}
    \textbf{Four abstraction levels:}

    \vspace{0.15cm}
    \textcolor{errorstroke}{\textbf{L1 Business:}} Revenue, engagement\\
    {\scriptsize Few metrics, high signal, exec attention}

    \vspace{0.1cm}
    \textcolor{routingstroke}{\textbf{L2 Portfolio:}} RecSys lift, fraud rate\\
    {\scriptsize Domain-level aggregation}

    \vspace{0.1cm}
    \textcolor{computestroke}{\textbf{L3 Model:}} Accuracy, latency, drift\\
    {\scriptsize Investigation on anomaly only}

    \vspace{0.1cm}
    \textcolor{datastroke}{\textbf{L4 Infra:}} GPU util, feature store\\
    {\scriptsize Platform team routing}

    \vspace{0.1cm}
    \mlsysconcept{Principle:}{Alert upward, investigate downward.}
  \end{column}
\end{columns}

\end{frame}

\begin{frame}{Drift Detection: PSI}
\note{[2 min] Two types of drift: covariate (input $P(X)$ changes) and
concept ($P(Y|X)$ changes). PSI is cheap to compute and serves as a leading
indicator. Thresholds: $<$0.1 stable, 0.1--0.25 investigate, $\geq$0.25 act.
Monitoring lag: 2\% accuracy drop needs $\sim$2,500 labeled samples.}

\footnotesize
\textbf{Population Stability Index} detects feature distribution shift:

\vspace{0.05cm}
$$PSI = \sum_{i=1}^{n} (A_i - E_i) \times \ln\left(\frac{A_i}{E_i}\right)$$

\vspace{0.05cm}
\renewcommand{\arraystretch}{1.1}
\begin{tabular}{@{}lll@{}}
  \toprule
  \textbf{PSI Value} & \textbf{Interpretation} & \textbf{Action} \\
  \midrule
  $< 0.1$          & No significant shift   & Continue monitoring \\
  $0.1$ -- $0.25$  & Moderate shift          & Investigate \\
  $\geq 0.25$      & Significant shift       & \textcolor{errorstroke}{\textbf{Immediate action}} \\
  \bottomrule
\end{tabular}

\vspace{0.1cm}
\mlsysalert{Monitoring lag:}{Detecting a 2\% accuracy drop needs $\sim$2,500 samples. At 1K labels/hr $\to$ \textbf{2.5 hrs} broken. PSI on inputs is faster.}

\end{frame}

% --- ACTIVE LEARNING 3: Discussion ---
\begin{frame}{Discussion: Which Monitoring Level Fires First?}
\note{[3 min] Turn-and-talk. Students discuss in pairs for 90 seconds.
The answer depends on the failure mode: infrastructure failures show at L4 first,
while data drift shows at L3 then propagates to L2/L1.}

\centering
\vspace{0.8cm}
{\large\bfseries Turn and Talk \textcolor{midgray}{(90 seconds)}}

\vspace{0.6cm}
{\large An upstream feature pipeline silently starts\\
returning null values for 3\% of users.\\[0.3cm]
\alert{Which monitoring level detects this first?\\
How long until the business metric dashboard shows a dip?}}

\vspace{0.6cm}
\begin{columns}[c]
  \begin{column}{0.22\textwidth}\centering\small L1\\Business\end{column}
  \begin{column}{0.22\textwidth}\centering\small L2\\Portfolio\end{column}
  \begin{column}{0.22\textwidth}\centering\small L3\\Model\end{column}
  \begin{column}{0.22\textwidth}\centering\small L4\\Infra\end{column}
\end{columns}

\end{frame}

% =============================================================================
\section{Features \& Organization}
% =============================================================================

\begin{frame}{Feature Freshness: Batch vs.\ Streaming}
\note{[3 min] Feature freshness latency determines recommendation quality.
Batch: 12--24 hours. Streaming: 1--5 seconds. The 10,000$\times$ improvement
yields 10--20\% engagement lift. Ask: ``For fraud detection, how stale can
features be before adversaries exploit the window?''}

% --- Full-width diagram ---
\centering
\safeimg[width=0.88\textwidth,height=4.5cm,keepaspectratio]{feature-freshness.pdf}

\vspace{0.1cm}
\small
$L_{\text{freshness}} = T_{\text{available}} - T_{\text{event}}$

\vspace{0.1cm}
{\footnotesize Streaming: \textbf{10--20\% engagement lift}. For fraud detection, day-old features give adversaries a 24-hour exploitation window.}

\end{frame}

\begin{frame}{Organizational Patterns for ML Platforms}
\note{[2 min] Three patterns: centralized (consistency but bottleneck),
embedded (velocity but fragmentation), hybrid (balanced). Most mature
organizations adopt hybrid. If short: just show the comparison table.}

\footnotesize
\renewcommand{\arraystretch}{1.15}
\begin{tabular}{@{}lp{3.2cm}p{3.2cm}p{3.2cm}@{}}
  \toprule
  & \textbf{Centralized} & \textbf{Embedded} & \textbf{Hybrid} \\
  \midrule
  \textbf{Structure}  & Single platform team & ML eng in each team & Core + domain leads \\
  \textbf{Strength}   & Consistency          & Responsiveness       & Balanced \\
  \textbf{Risk}       & Bottleneck           & Fragmentation        & Coordination cost \\
  \textbf{Best for}   & 100+ models          & 10--20 models        & 50--100 models \\
  \bottomrule
\end{tabular}

\vspace{0.15cm}
\begin{mlsyscard}{computestroke}
{\footnotesize \textbf{Decision factors:} Model count, model similarity, org size, regulatory requirements, and infrastructure maturity determine the optimal pattern.}
\end{mlsyscard}

\end{frame}

\begin{frame}{TCO Framework for ML Systems}
\note{[3 min] The four-component cost structure evolves with scale. Early-stage:
iteration-dominated (optimize velocity). Production: inference-dominated
(optimize efficiency). Walk through the equation and the shift.
Ask: ``Which cost component dominates at your organization?''}

\small
$$\TCO = C_{\text{train}} + C_{\text{infer}} + C_{\text{data}} + C_{\text{iter}}$$

\vspace{0.15cm}
{\footnotesize
\renewcommand{\arraystretch}{1.2}
\begin{tabular}{@{}lll@{}}
  \toprule
  \textbf{Component} & \textbf{Early Stage} & \textbf{Production Scale} \\
  \midrule
  $C_{\text{train}}$  & Dominant (experimentation)     & Fixed (scheduled retraining) \\
  $C_{\text{infer}}$  & Negligible (few users)         & \textcolor{errorstroke}{\textbf{Dominant}} (70--80\% of TCO) \\
  $C_{\text{data}}$   & Small (curated datasets)       & Growing (storage + pipelines) \\
  $C_{\text{iter}}$   & \textcolor{routingstroke}{\textbf{Dominant}} (engineer time)  & Reduced by automation \\
  \bottomrule
\end{tabular}
}

\vspace{0.15cm}
\mlsysconcept{Strategy:}{Match optimization investment to current cost structure.}

\end{frame}

% =============================================================================
\section{Wrap-Up}
% =============================================================================

\begin{frame}{Fallacies}
\note{[2 min] Four common misconceptions with quantitative evidence.}

\footnotesize
\textbf{Fallacy:} \textit{Operational complexity grows linearly with model count.}\\
100 models $\times$ 10 metrics at 5\% FPR $\to$ 14,400 false alerts/day. Dependencies grow as $O(N^2)$.

\vspace{0.08cm}
\textbf{Fallacy:} \textit{Platform investment only makes sense after 100+ models.}\\
A \$2M platform breaks even at $\sim$20 models. By 100, migration cost is 3--5$\times$ the initial build.

\vspace{0.08cm}
\textbf{Fallacy:} \textit{Monitoring training metrics provides sufficient observability.}\\
20 ms upstream slowdown $\to$ 30\% latency degradation across 10--50 downstream models.

\vspace{0.08cm}
\textbf{Fallacy:} \textit{Technical debt is inevitable and should wait.}\\
Monitoring debt: \$750K/yr (15 incidents $\times$ \$50K). Teams see 70--80\% capacity on maintenance.

\end{frame}

\begin{frame}{Pitfalls}
\note{[2 min] Three operational pitfalls. If short: cover just the first two.}

\footnotesize
\textbf{Pitfall:} \textit{Applying independent alerting to every model and metric.}\\
With $N=1000$ independent tests at $\alpha=0.05$: $P(\text{false alert}) = 1 - 0.95^{1000} \approx 1.0$. Operators ignore the noise; genuine incidents vanish. Use \textbf{hierarchical monitoring}.

\vspace{0.15cm}
\textbf{Pitfall:} \textit{Treating all deployments with uniform rollout procedures.}\\
Fraud detection needs hourly updates with seconds-fast rollback. LLMs need multi-day shadow testing. \textbf{Risk-based policies} must match model type.

\vspace{0.15cm}
\textbf{Pitfall:} \textit{Defaulting to batch feature computation for all features.}\\
Daily batch yields $L_{\text{freshness}} \approx 12$--24 hrs; streaming achieves 1--5 seconds. For recommendations, streaming yields \textbf{10--20\% engagement lift}; for fraud, batch gives adversaries a 24-hour window.

\end{frame}

% --- RETRIEVAL PRACTICE ---

% --- MUDDIEST POINT ---
\begin{frame}{Muddiest Point}
\note{[2 min] Quick anonymous poll. Students write on a slip of paper or submit
digitally. Collect and scan for patterns. Address the top 2--3 confusions in the
next lecture's opening. This closes the feedback loop.}

\centering
\vspace{1.0cm}
{\Large\bfseries What was the \alert{muddiest point} today?}

\vspace{0.8cm}
{\normalsize Write down the concept you found \textbf{most confusing}.}

\vspace{0.5cm}
{\small\textcolor{midgray}{Anonymous. One sentence. Submit before you leave.}}

\end{frame}

\begin{frame}{What Were the Key Ideas?}
\note{[2 min] Retrieval practice. Students write 90 seconds, no notes.
Do NOT show next slide yet. Walk around the room.}

\centering
\vspace{1.5cm}
{\Large\bfseries Close your notes.}

\vspace{0.8cm}
{\large Write down the \textbf{4 most important concepts} from today.}

\vspace{0.8cm}
{\normalsize\textcolor{midgray}{90 seconds --- no peeking.}}

\end{frame}

\begin{frame}{Key Takeaways}
\note{[2 min] Reveal. Walk through each bullet. Emphasize quantitative
anchors: 56\% savings, 14,400 alerts, 20$\times$ risk reduction,
\$1.08M daily loss.}

\scriptsize
\begin{itemize}\setlength\itemsep{0pt}
  \item \textbf{Platform ROI is superlinear}: Break-even at 20--50 models. 50 models saves 56\% (\$2M/yr).
  \item \textbf{Dependencies grow as $O(N^2)$}: Per-model practices collapse at $\sim$50 models.
  \item \textbf{Risk profile determines cadence}: LLMs: months. Fraud: hours. One-size CI/CD fails.
  \item \textbf{Alert fatigue is mathematical}: 100 models $\times$ 10 metrics $\to$ 14,400 false alerts/day. Use hierarchical monitoring.
  \item \textbf{Staged rollouts reduce risk 20$\times$}: 5\% canary limits blast radius.
  \item \textbf{TCO shifts with scale}: Early = iteration-dominated; production = inference-dominated (70--80\%).
  \item \textbf{Feature freshness matters}: Streaming yields 10--20\% engagement lift over batch.
\end{itemize}

\end{frame}

\begin{frame}{References}
\note{[1 min] Point students to canonical papers.}

\small
\mlsysref{Sculley+15}{Sculley et al. ``Hidden Technical Debt in ML Systems.'' NeurIPS 2015.}
\mlsysref{Covington+16}{Covington et al. ``Deep Neural Networks for YouTube Recommendations.'' RecSys 2016.}
\mlsysref{Chapelle+12}{Chapelle et al. ``Large-Scale Validation and Analysis of Interleaved Search Evaluation.'' TOIS 2012.}
\mlsysref{Kohavi+09}{Kohavi et al. ``Controlled Experiments on the Web.'' KDD 2009.}
\mlsysref{Hardt+16}{Hardt et al. ``Equality of Opportunity in Supervised Learning.'' NeurIPS 2016.}
\mlsysref{Paleyes+22}{Paleyes et al. ``Challenges in Deploying ML.'' ACM Comp.\ Surveys, 2022.}

\end{frame}

\begin{frame}{Next Lecture: Security and Privacy}
\note{[1 min] Forward hook. A globally distributed fleet that autonomously
adapts presents an expansive attack surface. Data poisoning, model extraction,
and differential privacy are the next challenges.}

\centering
\vspace{0.8cm}
{\large A fleet that is globally distributed\\
and autonomously adapting presents an\\
\alert{expansive attack surface.}}

\vspace{0.6cm}
\begin{columns}[c]
  \begin{column}{0.30\textwidth}
    \centering
    \begin{mlsyscard}{errorstroke}
    {\small\bfseries Data Poisoning}\\[0.1cm]
    {\scriptsize Corrupt training data}
    \end{mlsyscard}
  \end{column}
  \begin{column}{0.30\textwidth}
    \centering
    \begin{mlsyscard}{computestroke}
    {\small\bfseries Model Extraction}\\[0.1cm]
    {\scriptsize Steal model via API}
    \end{mlsyscard}
  \end{column}
  \begin{column}{0.30\textwidth}
    \centering
    \begin{mlsyscard}{datastroke}
    {\small\bfseries Differential Privacy}\\[0.1cm]
    {\scriptsize Protect training data}
    \end{mlsyscard}
  \end{column}
\end{columns}

\vspace{0.4cm}
{\normalsize\textbf{How do we defend the fleet?}}

\end{frame}



\appendix

\begin{frame}{Backup: Extended Reference}
\note{Backup slide with additional reference material for this chapter.}

\footnotesize
This slide provides extended reference material for students who want to go deeper.
Refer to the textbook chapter for complete derivations and additional worked examples.

\vspace{0.3cm}
\begin{mlsyscard}{crimson}
\textbf{Key equations and frameworks from this chapter} are collected in the
textbook's summary tables. Use them as a quick reference during problem sets.
\end{mlsyscard}

\end{frame}

\begin{frame}{Backup: Further Reading}
\note{Backup slide. Point students to additional resources beyond the references slide.}

\footnotesize
\textbf{For deeper exploration:}

\vspace{0.15cm}
\begin{itemize}\setlength\itemsep{2pt}
  \item The textbook chapter contains 2--3 additional worked examples
  \item Problem sets apply these concepts to real-world scenarios
  \item The course website links to open-source implementations
  \item Office hours: bring your toughest questions
\end{itemize}

\end{frame}


\end{document}
